\documentclass[11pt,a4paper]{scrartcl}

\usepackage[lithuanian]{babel}
\usepackage{../../1.\space Vienanariai\space ir\space daugianariai/manostilius_lt}
\usepackage{tasks}
\renewcommand{\theproblem}{\arabic{problem}}

\title{Kvadratinių lygčių uždaviniai}
\author{Andrius Berniukevičius}

\begin{document}

\maketitle
\section*{Kaip dirbti su uždaviniais?}

Pirmiausia kiekvieną uždavinį pabandykite išspręsti savarankiškai. Skirkite laiko apgalvoti, ką jau žinote ir kokį sprendimo būdą galėtumėte pritaikyti.

Paties pagimdytas sprendimas yra šimtą kartų vertingesnis už įvaikintą sprendimą.

\section{Uždaviniai}

\begin{problem}
    Įrodykite, kad su bet kokiomis realiosiomis parametrų $a$, $b$ ir $c$ reikšmėmis lygtis
    \[
    x^2+(a+b)x+ab-c^2=0
    \]
    visada turi bent vieną realųjį sprendinį. Kada ši lygtis turi lygiai vieną sprendinį?
\end{problem}

\begin{problem}
    Tegul $x_1$ ir $x_2$ yra lygties $x^2+px+q=0$ sprendiniai. Naudodamiesi Vijeto teorema įrodykite tapatybę:
    \[
    x_1^4+x_2^4=p^4-4p^2q+2q^2.
    \]
\end{problem}

\begin{problem}
    Duotos dvi skirtingos kvadratinės lygtys:
    \[
    x^2+px+q=0
    \qquad\text{ir}\qquad
    x^2+qx+p=0,
    \]
    kur $p\ne q$. Įrodykite, kad jeigu šios lygtys turi bendrą realųjį sprendinį, tai $p+q+1=0$.
\end{problem}

\begin{problem}
    Žinoma, kad kvadratinė lygtis $ax^2+bx+c=0$, kur $a\ne0$, neturi realiųjų sprendinių, o $a>0$. Įrodykite nelygybę
    \[
    a+b+c>0.
\]
\end{problem}

\begin{problem}
    Duotos dvi kvadratinės lygtys:
    \[
    x^2+ax+b=0
    \qquad\text{ir}\qquad
    x^2+cx+d=0.
    \]
    Jų koeficientai susieti lygybe $ac=2(b+d)$. Įrodykite, kad bent viena iš šių lygčių turi realiųjų sprendinių.
\end{problem}

\begin{problem}
    Kvadratinės lygties $ax^2+bx+c=0$ sprendiniai yra $x_1$ ir $x_2$. Įrodykite, kad jeigu vienas sprendinys yra $n$ kartų didesnis už kitą, t. y. $x_2=nx_1$, tai koeficientai tenkina lygybę
    \[
    nb^2=(n+1)^2ac.
    \]
\end{problem}

\begin{problem}
    Įrodykite, kad jeigu $a$, $b$ ir $c$ yra nelyginiai sveikieji skaičiai, tai lygtis $ax^2+bx+c=0$ negali turėti sveikojo sprendinio.
\end{problem}

\begin{problem}
    Įrodykite, kad jeigu $p$ ir $q$ yra sveikieji skaičiai, o lygties $x^2+px+q=0$ diskriminantas yra sveikojo skaičiaus kvadratas, tai visi lygties sprendiniai yra sveikieji skaičiai.
\end{problem}

\begin{problem}
    Įrodykite, kad jeigu lygtis $x^2+px+q=0$ turi realiųjų sprendinių, tai ir lygtis
    \[
    x^2+px+q+\left(x+\frac{p}{2}\right)^2=0
    \]
    būtinai turės realiųjų sprendinių.
\end{problem}

\begin{problem}
    Žinoma, kad lygties $x^2-px+q=0$ sprendiniai $x_1$ ir $x_2$ yra teigiami realieji skaičiai. Iš diskriminanto savybių išveskite ir įrodykite, kad $p\ge 2\sqrt q$.
\end{problem}

\end{document}
